% -------------------------------------------------------------------------------------
\section{领头阶共线拉氏量与共线 Wilson 线}\label{sec:scet-lagrangian}

令 $\xi_n\equiv P_n\psi$、$\phi_{\bar n}\equiv P_{\bar n}\psi$，于是 $\sla n\,\xi_n=0$，$\bar\xi_n=\bar\psi P_{\bar n}$，$\psi=\xi_n+\phi_{\bar n}$。对 $n$ 共线夸克，$\xi_n$ 是“好”分量，$\phi_{\bar n}$ 是被 $\lambda$ 压低的“坏”分量；SCET 的共线拉氏量就是在树图层次上把后者积掉的结果。

\begin{proposition}[领头阶共线夸克拉氏量 \tagR]\label{prop:scet-lagrangian}
（$R_{\mathrm{PT}}$：在重整化微扰论中对红外正规化的部分子态严格。）（树图层次的匹配。）
(a) 在经典作用量 $\int\bar\psi\,\ii\sla D\,\psi$（\cref{def:prelim-qcd}）中解出 $\phi_{\bar n}$ 的运动方程并代回，精确地得到
\begin{equation}\label{eq:scet-Lxixi-exact}
\mathcal L_{\xi\xi}=\bar\xi_n\Big[\ii n\cdot D+\ii\sla D_\perp\,\frac{1}{\ii\bar n\cdot D}\,\ii\sla D_\perp\Big]\frac{\bar{\sla n}}{2}\,\xi_n .
\end{equation}
(b) 令 $A=A_n+A_{us}$，超软场按\cref{eq:scet-multipole}作多极展开，并按\cref{prop:scet-scalings}只保留领头幂次，得
\begin{equation}\label{eq:scet-L0xi}
\mathcal L^{(0)}_{\xi\xi}=\bar\xi_n\Big[\ii n\cdot D+\ii\sla D_{\perp c}\,\frac{1}{\ii\bar n\cdot D_c}\,\ii\sla D_{\perp c}\Big]\frac{\bar{\sla n}}{2}\,\xi_n,\qquad
\begin{aligned}
&\ii n\cdot D=\ii n\cdot\partial+g\,n\cdot A_n+g\,n\cdot A_{us}(x_-),\\
&\ii\bar n\cdot D_c=\ii\bar n\cdot\partial+g\,\bar n\cdot A_n,\quad \ii D^\mu_{\perp c}=\ii\partial^\mu_\perp+gA^\mu_{n\perp}.
\end{aligned}
\end{equation}
% VF-SCET-3
(c) 相应的 Feynman 规则为：共线夸克传播子 $\ii\frac{\sla n}{2}\frac{\bar n\cdot p}{p^2+\ii0}$；超软胶子顶角 $\ii g\,t^a\,n^\mu\frac{\bar{\sla n}}{2}$；共线胶子顶角（夸克动量 $p\to p'$，胶子动量 $p'-p$ 流入）
\begin{equation}\label{eq:scet-collvertex}
V^\mu=\ii g\,t^a\Big[n^\mu+\frac{\gamma_\perp^\mu\sla p_\perp}{\bar n\cdot p}+\frac{\sla p'_\perp\gamma_\perp^\mu}{\bar n\cdot p'}-\frac{\sla p'_\perp\sla p_\perp}{\bar n\cdot p\,\bar n\cdot p'}\,\bar n^\mu\Big]\frac{\bar{\sla n}}{2},
\end{equation}
% VF-SCET-4
其中 $\gamma_\perp^\mu=\gamma^\mu-n^\mu\frac{\bar{\sla n}}{2}-\bar n^\mu\frac{\sla n}{2}$；并且对任意离壳动量精确成立
$P_n\,S(p')\,(\ii g\gamma^\mu t^a)\,S(p)\,P_{\bar n}=S_\xi(p')\,V^\mu\,S_\xi(p)$，$S(p)=\frac{\ii\sla p}{p^2+\ii0}$，$S_\xi$ 为上述共线传播子。（本命题只涉及部分子场与部分子振幅。）
\uses{prop:scet-scalings, def:prelim-qcd}
\end{proposition}

\begin{proof}
(a) 把 $\sla D=\frac{\sla n}{2}\bar n\cdot D+\frac{\bar{\sla n}}{2}n\cdot D+\sla D_\perp$ 代入，并用约定段落中投影算符的性质（$\sla n^2=0$、$\bar{\sla n}P_n=\bar{\sla n}$、$\gamma_\perp$ 与 $P_n$ 对易）得四个分块恒等式
$P_{\bar n}\sla DP_n=n\cdot D\,\frac{\bar{\sla n}}{2}P_n$、$P_n\sla DP_{\bar n}=\bar n\cdot D\,\frac{\sla n}{2}P_{\bar n}$、$P_{\bar n}\sla DP_{\bar n}=\sla D_\perp P_{\bar n}$、$P_n\sla DP_n=\sla D_\perp P_n$。由 $\bar\xi_n=\bar\psi P_{\bar n}$、$\bar\phi_{\bar n}=\bar\psi P_n$，
\begin{equation}\label{eq:scet-psiDpsi}
\bar\psi\,\ii\sla D\,\psi=\bar\xi_n\frac{\bar{\sla n}}{2}\,\ii n\cdot D\,\xi_n+\bar\xi_n\,\ii\sla D_\perp\phi_{\bar n}
+\bar\phi_{\bar n}\,\ii\sla D_\perp\xi_n+\bar\phi_{\bar n}\frac{\sla n}{2}\,\ii\bar n\cdot D\,\phi_{\bar n}.
\end{equation}
对 $\bar\phi_{\bar n}$ 变分得 $\frac{\sla n}{2}\ii\bar n\cdot D\,\phi_{\bar n}+\ii\sla D_\perp\xi_n=0$；左乘 $\frac{\bar{\sla n}}{2}$ 并用 $\frac{\bar{\sla n}\sla n}{4}\phi_{\bar n}=\phi_{\bar n}$ 得 $\phi_{\bar n}=-\frac{1}{\ii\bar n\cdot D}\frac{\bar{\sla n}}{2}\,\ii\sla D_\perp\xi_n$。代回后含 $\bar\phi_{\bar n}$ 的两项之和按运动方程为零，剩下 $\bar\xi_n\frac{\bar{\sla n}}2\ii n\cdot D\,\xi_n-\bar\xi_n\ii\sla D_\perp\frac{1}{\ii\bar n\cdot D}\frac{\bar{\sla n}}{2}\ii\sla D_\perp\xi_n$；把 $\bar{\sla n}$ 越过 $\sla D_\perp$ 移到最右边（反对易，多一个负号）即得\cref{eq:scet-Lxixi-exact}。这一步没有任何近似；在路径积分中对 $\phi_{\bar n}$ 作 Grassmann--Gauss 积分还给出一个与 $\xi_n$ 无关的泛函行列式，它不影响树图。

(b) 由\cref{prop:scet-scalings}，作用在共线场上时 $\ii n\cdot\partial$、$g\,n\cdot A_n$、$g\,n\cdot A_{us}$ 都是 $\order{\lambda^2}$，故全部保留，后者按\cref{eq:scet-multipole}取 $A_{us}(x_-)$；$\ii\bar n\cdot\partial$ 与 $g\,\bar n\cdot A_n$ 是 $\order{1}$，$g\,\bar n\cdot A_{us}\sim\lambda^2$ 被压低；$\ii\partial_\perp$ 与 $gA_{n\perp}$ 是 $\order{\lambda}$，$gA_{us\perp}\sim\lambda^2$ 被压低。代入\cref{eq:scet-Lxixi-exact}即得\cref{eq:scet-L0xi}。

(c) 自由部分：方括号作用在 $P_n$ 子空间上等于 $n\cdot p+\sla p_\perp\sla p_\perp/\bar n\cdot p=p^2/\bar n\cdot p$，其逆乘以 $\ii$ 即传播子，因 $\ii\frac{\sla n}{2}\frac{\bar n\cdot p}{p^2}\cdot\frac{p^2}{\bar n\cdot p}\frac{\bar{\sla n}}{2}=\ii P_n$；它等于\cref{eq:scet-xiprop}。超软顶角来自 $\bar\xi_n\,g\,n\cdot A_{us}\frac{\bar{\sla n}}2\xi_n$。共线顶角是\cref{eq:scet-L0xi}在 $A_{us}=0$ 时的 $\order{g}$ 项：$n^\mu$ 来自 $\ii n\cdot D$，中间两项来自两个 $\ii\sla D_{\perp c}$ 中的 $gA_{n\perp}$，末项来自 $\frac{1}{\bar n\cdot p+g\bar n\cdot A_n}=\frac{1}{\bar n\cdot p}-\frac{1}{\bar n\cdot p'}g\,\bar n\cdot A_n\frac{1}{\bar n\cdot p}+\order{g^2}$。最后，(a) 表明消去 $\phi_{\bar n}$ 后的树图生成泛函等于 QCD 的树图生成泛函对 $\xi_n$ 源的限制；而对共线胶子，(b) 没有丢掉任何项，所以 $\xi_n$ 的单胶子树图关联函数 $S_\xi V^\mu S_\xi$ 精确地等于 QCD 关联函数的投影 $P_n S\,(\ii g\gamma^\mu t^a)\,S\,P_{\bar n}$（对随机的离壳动量，两边逐元素之差的相对大小 $<10^{-15}$）。% w18_s1 (D)
\end{proof}

\paragraph{共线 Wilson 线。}
由\cref{prop:scet-scalings}，$g\,\bar n\cdot A_n\sim\lambda^0$，故必须把它对所有阶求和。定义（\cref{def:prelim-wilson}的约定，离上端点近者在左）
\begin{equation}\label{eq:scet-Wn}
W_n(x)\equiv W[x,x-\infty\bar n]=\mathrm P\exp\Big(\ii g\int_{-\infty}^0\dd s\;\bar n\cdot A_n(x+s\bar n)\Big)
=\sum_{\rm perms}\exp\Big[-\frac{g}{\bar{\mathcal P}}\,\bar n\cdot A_n\Big],
\end{equation}
% VF-SCET-5
最后一式是 BFPS 的标签动量形式\cite{BauerEtAl2001}。由\cref{prop:prelim-wilson-covariance}，$W_n^\dagger W_n=1$，$\bar n\cdot\partial\,W_n=\ii g\,\bar n\cdot A_n\,W_n$，即 $\bar n\cdot D_c\,W_n=0$，从而
$W_n^\dagger\,\ii\bar n\cdot D_c\,W_n=\ii\bar n\cdot\partial$、$\frac{1}{\ii\bar n\cdot D_c}=W_n\frac{1}{\ii\bar n\cdot\partial}W_n^\dagger$；在 $U_c(x-\infty\bar n)=1$ 的共线规范变换下 $W_n\to U_c(x)W_n$，所以 $\chi_n\equiv W_n^\dagger\xi_n$ 是共线规范不变的。单胶子规则：被吸收的胶子（场 $\propto\epsilon\,\ee^{-\ii k\cdot y}$）给出 $\ii g\,\bar n\cdot\epsilon\,t^a\int_{-\infty}^0\dd s\,\ee^{-\ii s\bar n\cdot k+0s}=-g\,t^a\,\bar n\cdot\epsilon/(\bar n\cdot k+\ii0)$；产生的出射胶子给出 $+g\,t^a\,\bar n\cdot\epsilon^*/(\bar n\cdot k-\ii0)$；两者都是\cref{prop:prelim-wilson-rules}中指向过去的线的规则（$n\to\bar n$）。

\begin{proposition}[共线 Wilson 线 \tagR]\label{prop:scet-wilson}
（$R_{\mathrm{PT}}$：在重整化微扰论中对红外正规化的部分子态严格。）（树图层次。）设 QCD 矢量流 $J^\mu=\bar\psi\gamma^\mu\psi$ 产生一个 $n$ 共线夸克、一个 $\bar n$ 共线反夸克及任意个 $n$ 共线与 $\bar n$ 共线胶子。积掉虚度 $\sim Q^2$ 的离壳线之后，其树图矩阵元在 $\lambda$ 的领头幂次等于
\begin{equation}\label{eq:scet-Jtree}
J^\mu_{\rm SCET}=\bar\xi_n\,W_n\,\gamma_\perp^\mu\,W_{\bar n}^\dagger\,\xi_{\bar n}=\bar\chi_n\gamma_\perp^\mu\chi_{\bar n}
\end{equation}
用 $\mathcal L^{(0)}$ 算出的矩阵元，其中 $W_n$ 由\cref{eq:scet-Wn}给出，单胶子规则为 $-g\,\bar n^\mu t^a/(\bar n\cdot k+\ii0)$。与 BFPS、BPS 的公式相比不需要改变 $g$ 的符号。（本命题只涉及部分子振幅。）
\uses{prop:scet-lagrangian, def:prelim-wilson, prop:prelim-wilson-covariance, prop:prelim-wilson-rules, lem:prelim-eikonal-identity}
\end{proposition}

\begin{proof}
考虑末态为 $q(p_1)\bar q(p_2)g(k_1)\cdots g(k_m)$ 的任一树图，先设胶子都是 $n$ 共线的。\emph{第一步（分解）。}从流顶点到外反夸克的线携带 $p_2+K$，$K$ 是若干 $n$ 共线动量之和；由 $2p_2\cdot K=n\cdot p_2\,\bar n\cdot K+\bar n\cdot p_2\,n\cdot K+2p_{2\perp}\cdot K_\perp$ 与\cref{def:scet-modes}，$(p_2+K)^2=n\cdot p_2\,\bar n\cdot K\,[1+\order{\lambda^2}]\sim Q^2$，这些线是要积掉的离壳线。接在这条线上的每个胶子都是一棵只含 $n$ 共线线（虚度 $\sim\lambda^2Q^2$）的子树之根；子树、外夸克线以及接在 $n$ 共线夸克线上的胶子都属于共线扇区，由\cref{prop:scet-lagrangian}在树图层次精确描述。\emph{第二步（领头幂次）。}设反夸克线上依次（自外反夸克向流顶点）接了 $r$ 棵子树，其流为 $\mathcal J_j$、色指标 $a_j$，记 $K_j$ 为前 $j$ 棵子树的总动量。反夸克线给出
\[
\Gamma^\mu\,\frac{\ii(-\sla p_2-\sla K_r)}{(p_2+K_r)^2+\ii0}\,(\ii g\,\sla{\mathcal J}_r t^{a_r})\cdots\frac{\ii(-\sla p_2-\sla K_1)}{(p_2+K_1)^2+\ii0}\,(\ii g\,\sla{\mathcal J}_1t^{a_1})\,v(p_2).
\]
在领头幂次 $-\sla p_2-\sla K_j\to-\frac{n\cdot p_2}{2}\bar{\sla n}$；$\sla{\mathcal J}_j=\frac{\bar n\cdot\mathcal J_j}{2}\sla n+\frac{n\cdot\mathcal J_j}{2}\bar{\sla n}+\sla{\mathcal J}_{j\perp}$ 中，按\cref{prop:scet-scalings}后两项分别为 $\order{\lambda^2}$ 与 $\order{\lambda}$（且 $\bar{\sla n}$ 项紧挨 $\bar{\sla n}$ 而为零），于是每一对“传播子 × 顶角”化为 $\frac{g\,\bar n\cdot\mathcal J_j}{4(\bar n\cdot K_j+\ii0)}\,\bar{\sla n}\sla n\,t^{a_j}$；由于 $\bar{\sla n}v(p_2)=\order{\lambda}$，$\bar{\sla n}\sla n\,v=(4-\sla n\bar{\sla n})v=4v+\order{\lambda}$，整条线化为
$\Gamma^\mu v(p_2)\prod_{j=r}^{1}\big[g\,t^{a_j}\,\bar n\cdot\mathcal J_j/(\bar n\cdot K_j+\ii0)\big]$（$j=r$ 在最左）。又 $\bar u_n\gamma^\mu v_{\bar n}=\bar u_n\gamma_\perp^\mu v_{\bar n}$（因 $\bar u_n\sla n=0$、$\bar{\sla n}v_{\bar n}=0$）。\emph{第三步（对接法求和）。}对固定的一组子树，对其根在线上的 $r!$ 种次序求和。另一方面，\cref{eq:scet-Wn}对 $r$ 个出射胶子的项在给定次序 $s_1<\cdots<s_r\le0$ 上的嵌套积分逐层初等可积，$\int_{-\infty}^{s_{j+1}}\dd s_j\,\ee^{\ii s_j\bar n\cdot K_j}\cdots=\frac{\ee^{\ii s_{j+1}\bar n\cdot K_j}}{\ii\bar n\cdot K_j}\cdots$，给出 $(\ii g)^r\prod_j(\ii\,\bar n\cdot K_j)^{-1}=g^r/\prod_j\bar n\cdot K_j$，色矩阵按 $t^{a_r}\cdots t^{a_1}$ 排列（离 $x$ 近者在左）——这正是上式在同一次序下的值，而对次序求和正是路径有序指数的展开（\cref{prop:prelim-wilson-rules}；交换群极限即\cref{lem:prelim-eikonal-identity}）。两者的 $\ii0$ 不同（QCD 线给出未来指向的 $+\ii0$），但对共线胶子 $\bar n\cdot K_j\sim Q\neq0$，$\ii0$ 无关；$W_n$ 取自 $-\infty$ 是约定，方向依赖只来自 $\bar n\cdot k\to0$ 的区域，在 SCET 中被零箱减除抵消\cite{RothsteinStewart2016}。\emph{第四步。}于是每个 QCD 树图都等于子树接在 $W_n$ 的 $r$ 胶子顶点上的 SCET 图；$\bar n$ 共线胶子接在 $n$ 共线夸克线上时同理给出 $W_{\bar n}^\dagger$。这就证明了\cref{eq:scet-Jtree}。推导只用到 $D=\partial-\ii gA$ 与 $W=\mathrm P\exp(+\ii g\int\dd z\cdot A)$，而 BFPS、BPS 采用相同的约定（$\ii D=\ii\partial+gA$），故符号一致。
\end{proof}

%%% BEGIN fig:scet-wilson
\begin{figure}[htbp]
\centering
\begin{tikzpicture}[fd,
  t2 n/.style={draw=blue!60!black},
  t2 nb/.style={draw=orange!80!black},
  t2 us/.style={draw=green!45!black},
  t2 lab/.style={font=\small},
  t2 eq/.style={font=\large}]
  % ======================= (a) origin of the collinear Wilson line W_n =======================
  \begin{scope}
    \node[opvertex] (O) at (0,0) {};
    \coordinate (Q) at (2.2,0.95); \coordinate (Qb) at (2.2,-0.95);
    \coordinate (G) at ($(O)!0.42!(Qb)$);
    \draw[quark,t2 n] (O) -- (Q);
    \draw[quark,t2 nb] (Qb) -- (G); \draw[quark=0.62,t2 nb] (G) -- (O);
    \draw[gluon,t2 n] (G) -- (2.25,0.22);
    \node[vertex] at (G) {};
    \node[t2 lab,anchor=west] at (Q) {$n$ 共线夸克};
    \node[t2 lab,anchor=west] at (Qb) {$\bar n$ 共线反夸克};
    \node[t2 lab,anchor=west] at (2.25,0.22) {$n$ 共线胶子};
    \node[font=\footnotesize,anchor=north east] at ($(O)!0.5!(G)+(0.04,-0.06)$) {$p^2\sim Q^2$};
    \node[t2 eq] at (5.4,0) {$\Longrightarrow$};
    \node[font=\footnotesize,align=center] at (5.4,-0.55) {积掉\\离壳模式};
    \begin{scope}[xshift=6.7cm]
      \node[opvertex] (O) at (0,0) {};
      \draw[quark,t2 n] (O) -- (2.2,0.95) coordinate (Q);
      % the nbar-collinear antiquark and, parallel to it, the eikonal line W_n (along nbar)
      \draw[quark,t2 nb] ($(O.-67)+(-23.36:2.35)$) -- (O.-67);
      \draw[wilson=0.72] ($(O.20)+(-23.36:2.0)$) coordinate (W) -- (O.20);
      \coordinate (Y) at ($(O.20)!0.6!(W)$);
      \draw[gluon,t2 n] (Y) -- ($(Y)+(1.05,0.52)$) coordinate (Ge);
      \node[t2 lab,anchor=west] at (Q) {$n$ 共线夸克};
      \node[t2 lab,anchor=west] at (Ge) {$\bar n\cdot A_n$};
      \node[t2 lab,anchor=west] at (W) {$W_n$};
    \end{scope}
    \node[font=\small] at (4.6,-1.75) {$\bar\xi_n\,\Gamma\,\xi_{\bar n}\;\longrightarrow\;\bar\xi_n W_n\,\Gamma\,W_{\bar n}^\dagger\xi_{\bar n}=\bar\chi_n\,\Gamma\,\chi_{\bar n},\qquad \chi_n\equiv W_n^\dagger\xi_n$};
    \node[font=\small,anchor=west] at (-0.5,-2.45) {(a)\ 共线 Wilson 线 $W_n$（沿 $\bar n$ 方向，由 $\bar n\cdot A_n$ 构成）的来源};
  \end{scope}
  % ======================= (b) ultrasoft decoupling =======================
  \begin{scope}[yshift=-5.2cm]
    \node[opvertex] (O) at (0,0) {};
    \coordinate (Q) at (2.9,1.0);
    \coordinate (U1) at ($(O)!0.36!(Q)$); \coordinate (U2) at ($(O)!0.7!(Q)$);
    \draw[quark=0.2,t2 n] (O) -- (U1); \draw[quark,t2 n] (U1) -- (U2); \draw[quark,t2 n] (U2) -- (Q);
    \draw[gluon,t2 us] (U1) -- ($(U1)+(-0.3,1.1)$);
    \draw[gluon,t2 us] (U2) -- ($(U2)+(0.05,1.1)$);
    \foreach \p in {U1,U2} \node[vertex] at (\p) {};
    \node[t2 lab,anchor=west] at (Q) {$\xi_n$};
    \node[font=\footnotesize,anchor=west,text=green!35!black] at ($(U2)+(0.12,0.75)$) {$n\cdot A_{us}$};
    \node[t2 eq] at (4.75,0.45) {$=$};
    \begin{scope}[xshift=6.1cm]
      \node[opvertex] (O) at (0,0) {};
      % the n-collinear quark and, parallel to it, the eikonal line Y_n (along n)
      \draw[quark,t2 n] (O.-26) -- ++(19.03:3.0) coordinate (Q0);
      \draw[wilson=0.86] (O.64) -- ++(19.03:2.95) coordinate (W);
      \coordinate (Y1) at ($(O.64)!0.36!(W)$); \coordinate (Y2) at ($(O.64)!0.68!(W)$);
      \draw[gluon,t2 us] (Y1) -- ($(Y1)+(-0.3,1.0)$);
      \draw[gluon,t2 us] (Y2) -- ($(Y2)+(0.05,1.0)$);
      \node[t2 lab,anchor=north west] at (Q0) {$\xi_n^{(0)}$};
      \node[t2 lab,anchor=south west] at (W) {$Y_n$};
    \end{scope}
    \node[font=\small] at (4.4,-1.0) {$\xi_n=Y_n\,\xi_n^{(0)},\quad A_n=Y_n\,A_n^{(0)}\,Y_n^\dagger$};
    \node[font=\small,anchor=west] at (-0.5,-1.75) {(b)\ 超软退耦：超软胶子只与沿 $n$ 方向、由 $n\cdot A_{us}$ 构成的 Wilson 线 $Y_n$ 耦合};
  \end{scope}
\end{tikzpicture}
\caption{SCET 中的两类 Wilson 线。(a)~$n$ 共线胶子（蓝色，领头分量 $\bar n\cdot A_n\sim\lambda^0$）附着在 $\bar n$ 共线（或硬）线上时，使该线离壳 $p^2\sim Q^2$；
把这些离壳模式积掉后，它们附着到从流顶点 $\otimes$ 出发、沿 $\bar n$ 方向（与 $\bar n$ 共线反夸克平行）的共线 Wilson 线 $W_n$（由 $\bar n\cdot A_n$ 构成）上，
流算符因而只通过规范不变组合 $\chi_n=W_n^\dagger\xi_n$ 出现（树图层次的匹配）。
(b)~领头阶共线拉氏量中超软胶子（绿色）只通过 $n\cdot A_{us}$ 与 $n$ 共线夸克耦合；场重定义 $\xi_n=Y_n\xi_n^{(0)}$、$A_n=Y_nA_n^{(0)}Y_n^\dagger$
（$n\cdot D_{us}Y_n=0$）使新场 $\xi_n^{(0)}$、$A_n^{(0)}$ 与超软胶子退耦，超软胶子全部附着在沿 $n$ 方向（与 $n$ 共线夸克平行）的 Wilson 线 $Y_n$ 上。
$Y_n$ 的端点（$-\infty$ 或 $+\infty$）与 $\ii0$ 规定应依粒子是入射还是出射选取（此处为出射夸克）。Wilson 线上的箭头与被替代的费米子线的费米子数流向一致。}
\label{fig:scet-wilson}
\end{figure}
%%% END fig:scet-wilson

\begin{example}[$W_n$ 的树图匹配]\label{ex:scet-wn-tree}
取 $\gamma^*\to q(p_1)\bar q(p_2)g(k)$，$k\parallel p_1\parallel n$（\cref{fig:scet-wilson}(a)）。胶子接在夸克上的图由\cref{eq:scet-collvertex}的共线顶角与共线传播子精确再现。胶子接在反夸克上的图为
$\bar u(p_1)\gamma^\mu\frac{\ii(-\sla p_2-\sla k)}{(p_2+k)^2+\ii0}(\ii g\,t^a\sla\epsilon^*)v(p_2)$，其中 $(p_2+k)^2\simeq n\cdot p_2\,\bar n\cdot k\sim Q^2$。按\cref{prop:scet-wilson}证明的第二步（$r=1$），它在领头幂次等于
\[
g\,t^a\,\frac{\bar n\cdot\epsilon^*}{\bar n\cdot k}\;\bar u(p_1)\gamma_\perp^\mu v(p_2),
\]
即 $W_n$ 的单胶子项。在一个随机选取的运动学点上，该图与此式之差的相对大小在 $\lambda=10^{-1},10^{-2},10^{-3}$ 时为 $4.6\times10^{-2}$、$4.7\times10^{-3}$、$4.8\times10^{-4}$，按 $\lambda$ 的一次方减小；完整的 QCD 振幅与 SCET 振幅（共线顶角项加 $W_n$ 项）之差的相对大小为 $9.1\times10^{-3}$、$8.5\times10^{-5}$、$8.5\times10^{-7}$。% w18_s2 (B)
取 $\epsilon^*\to k$ 时 QCD 两图之和为零（Ward 恒等式）；在 SCET 中，共线顶角项给出 $-g\,t^a\bar u\gamma_\perp^\mu v$，$W_n$ 项给出 $+g\,t^a\bar u\gamma_\perp^\mu v$，二者精确相消：$W_n$ 恢复了共线规范不变性。
\uses{prop:scet-wilson}
\end{example}
